\documentclass[11pt,a4paper]{article}

\usepackage[utf8]{inputenc}
\usepackage[T1]{fontenc}
\usepackage{lmodern}
\usepackage[margin=1in]{geometry}
\usepackage{graphicx}
\usepackage{xcolor}
\usepackage{hyperref}
\usepackage{titling}
\usepackage{titlesec}
\usepackage{enumitem}
\usepackage{listings}
\usepackage{booktabs}
\usepackage{longtable}
\usepackage{fancyhdr}
\usepackage{abstract}

% Colours
\definecolor{geosinavy}{HTML}{0B2A4A}
\definecolor{geositeal}{HTML}{108A92}
\definecolor{geosimuted}{HTML}{555F6D}
\definecolor{codebg}{HTML}{F4F1EA}

\hypersetup{
    colorlinks=true,
    linkcolor=geosinavy,
    urlcolor=geositeal,
    citecolor=geosinavy,
    pdftitle={GeoSI: The Beginning of Geospatial Superintelligence},
    pdfauthor={Aaron R}
}

\titleformat{\section}{\Large\bfseries\color{geosinavy}}{\thesection}{1em}{}
\titleformat{\subsection}{\large\bfseries\color{geositeal}}{\thesubsection}{1em}{}

\lstset{
  basicstyle=\ttfamily\small,
  backgroundcolor=\color{codebg},
  frame=single,
  rulecolor=\color{geosimuted},
  breaklines=true,
  columns=fullflexible,
  keepspaces=true,
  xleftmargin=6pt,
  xrightmargin=6pt,
  framesep=6pt,
}

\pagestyle{fancy}
\fancyhf{}
\fancyhead[L]{\color{geosimuted}\small GeoSI --- Geospatial Superintelligence}
\fancyhead[R]{\color{geosimuted}\small v2.0}
\fancyfoot[C]{\thepage}

\title{%
  \vspace{-1cm}%
  {\color{geosinavy}\Huge\bfseries GeoSI}\\[0.4em]
  {\color{geositeal}\Large The Beginning of Geospatial Superintelligence}\\[0.3em]
  {\color{geosimuted}\normalsize A Conversational Multi-Tool Engine for Professional GIS}
}
\author{Aaron R \\ \small Digital University Kerala (DUK) \\ \small \href{mailto:aaronr.ds25@duk.ac.in}{aaronr.ds25@duk.ac.in}}
\date{May 2026 \quad $\cdot$ \quad Version 2.0}

\begin{document}

\maketitle
\thispagestyle{empty}

\begin{abstract}
\noindent
Professional Geographic Information Systems (GIS) offer enormous analytical
power but remain gated behind menus, scripts, and domain-specific query
languages. This report introduces \textbf{GeoSI} (Geospatial
Superintelligence), a conversational AI system that bridges natural
language and professional GIS. GeoSI decomposes plain-English spatial
questions into multi-step workflows and executes them across 125
geospatial tools, operating on layers already loaded in the user's QGIS
workspace. The system is built on a three-surface architecture that
strictly separates the universal \emph{engine} from the QGIS
\emph{plugin} and an optional \emph{REST server}, with Ollama as the
default local language model. We describe the system design, the
backend-dispatch pattern that keeps the engine framework-free, the
LLM priority chain with graceful fallback, the parser/agent/executor
pipeline, and an evaluation on representative vector and raster
workflows. GeoSI is production-ready and, we argue, represents a first
step toward domain-native geospatial superintelligence.
\end{abstract}

\vspace{0.5em}
\noindent\textbf{Keywords:} geospatial AI, natural language interface,
QGIS, large language models, multi-step planning, tool use, Ollama.

\tableofcontents
\newpage

\section{Introduction}
Geographic Information Systems are indispensable across urban planning,
disaster response, agriculture, conservation, and public policy. Yet
the cognitive overhead remains high. Translating a spatial question
into a correct toolchain requires fluency in algorithms, coordinate
reference systems, and vendor-specific interfaces. Domain experts who
understand the \emph{question} rarely produce the \emph{answer} without
GIS analyst intermediaries.

General-purpose AI assistants do not solve this gap: they can describe
GIS operations but cannot actually execute them on the user's data,
and they cannot guarantee that the data they reference has been loaded
at all.

\subsection{Contributions}
This report presents GeoSI, whose contributions are:
\begin{itemize}[leftmargin=*]
  \item A \textbf{three-surface architecture} that isolates QGIS
  dependencies to a single plugin layer while keeping the analytical
  engine framework-free and importable anywhere.
  \item A \textbf{backend-dispatch pattern} that allows the same tool
  specifications to run under QGIS Processing, a GeoPandas backend, or
  return a clear unavailability error.
  \item An \textbf{Ollama-first LLM chain} with 1.5~s availability
  probing and graceful fallback through Gemini, Anthropic, OpenAI,
  and a deterministic keyword parser, delivering a fully offline
  default experience.
  \item A \textbf{layer-grounding mechanism} that refuses to
  hallucinate data: queries referencing unloaded layers receive fuzzy
  suggestions based on the QGIS Layers panel.
  \item A \textbf{unified tool catalogue} of 125 auto-discovered tools
  across nine GIS domains, exposed through one conversational surface.
\end{itemize}

\section{System Architecture}
\subsection{Three-Surface Model}
GeoSI is organised into three strictly independent surfaces:
\begin{description}[leftmargin=*]
  \item[\texttt{geosi\_engine/}] The universal core. Zero QGIS imports.
  Contains models, tool registry, parser, agent, executor,
  state manager, validation, and conversation memory.
  \item[\texttt{geosi\_plugin/}] The only code that imports
  \texttt{qgis}. Owns the chat dock, syncs the QGIS Layers panel into
  engine state, and installs a \texttt{"qgis"} backend into the engine
  at startup.
  \item[\texttt{geosi\_server/}] An optional FastAPI REST surface. It
  can auto-fetch OpenStreetMap data when the referenced feature name
  is recognisable and safe.
\end{description}

\subsection{Backend Dispatch Pattern}
A \texttt{BackendTool} carries a specification (parameters, category,
description) but never imports QGIS. Execution is delegated to a
backend function registered at runtime:
\begin{lstlisting}[language=Python]
from geosi_engine.base import register_backend
register_backend("qgis", _run_qgis_algorithm)
\end{lstlisting}
This keeps the engine exportable as a pure Python library while
allowing the plugin to inject the QGIS Processing framework at boot.

\subsection{Parser, Agent, Executor}
The engine follows a classical plan-and-execute pipeline:
\begin{enumerate}[leftmargin=*]
  \item \textbf{Parser} produces an \texttt{AnalysisRequest} capturing
  intent (e.g. \textsc{Proximity}, \textsc{Raster}), entities
  (\texttt{primary\_layer}, \texttt{secondary\_layer}), and parameters
  (\texttt{distance\_value}, \texttt{threshold}).
  \item \textbf{Agent} validates that referenced layers exist, then
  produces an \texttt{ExecutionPlan} containing ordered \texttt{ToolStep}
  instances with parameter bindings and dependency edges.
  \item \textbf{Executor} walks the plan, dispatches each step to the
  registered backend, handles partial failures, and returns a
  structured \texttt{ExecutionResult}.
\end{enumerate}

\section{Language Model Integration}
\subsection{Priority Chain}
GeoSI attempts providers in order of local-first preference:
\begin{enumerate}[leftmargin=*]
  \item \textbf{Ollama} (local, default). 1.5~s probe against
  \texttt{/api/tags}; 30~s generate timeout.
  \item \textbf{Google Gemini} when \texttt{GEMINI\_API\_KEY} is set.
  \item \textbf{Anthropic Claude} when \texttt{ANTHROPIC\_API\_KEY} is set.
  \item \textbf{OpenAI} when \texttt{OPENAI\_API\_KEY} is set.
  \item \textbf{Rule-based keyword parser} --- always available, drives
  all 125 tools offline.
\end{enumerate}
Every provider failure is logged and falls through. The deterministic
fallback guarantees GeoSI always produces an answer, even without any
LLM.

\subsection{Prompt Design}
The planner prompt exposes the tool catalogue as structured JSON,
the user's available layers, and the query text. The LLM is asked to
return strictly valid JSON conforming to the \texttt{ExecutionPlan}
schema. Responses are sanitised (fenced-block removal, tolerant
extraction) before instantiation; malformed responses trigger the
keyword fallback.

\section{Tool Catalogue}
GeoSI auto-discovers 125 tools via \texttt{pkgutil.walk\_packages} over
\texttt{geosi\_engine.tools}. Table~\ref{tab:tools} summarises the
coverage.

\begin{table}[ht]
\centering
\caption{GeoSI tool catalogue (125 tools across nine domains).}
\label{tab:tools}
\begin{tabular}{@{}lll@{}}
\toprule
\textbf{Domain} & \textbf{Count} & \textbf{Representative Tools} \\
\midrule
Vector & 40 & buffer, intersect, clip, dissolve, voronoi, joins \\
Raster & 25 & raster\_calc, NDVI, NDWI, zonal\_stats, reclassify \\
Terrain & 10 & slope, aspect, hillshade, contour, viewshed \\
Network & 8 & shortest\_path, service\_area, isochrone, OD matrix \\
Temporal & 6 & change\_detection, time\_series\_aggregate \\
AI/ML & 10 & kmeans, dbscan, Getis-Ord Gi*, Moran's I, IDW \\
Cartography & 10 & export\_geojson, export\_kml, geocode \\
Validation & 10 & fix\_geometries, topology\_check, duplicates \\
Portable & 2 & count\_features, load\_spatial\_data \\
\bottomrule
\end{tabular}
\end{table}

\section{Layer Grounding}
A distinctive property of GeoSI is its refusal to hallucinate data.
When a user references a layer that is not in the QGIS Layers panel,
the agent responds with a clear error:
\begin{lstlisting}
User: "Buffer hospitals by 500m"
(only "Schools" is loaded)
GeoSI: Layer 'hospitals' not found in the QGIS Layers panel.
       Did you mean: 'Schools' (50% match)?
\end{lstlisting}
Two mechanisms collaborate:
\begin{enumerate}[leftmargin=*]
  \item A fuzzy matcher (SequenceMatcher) with a 0.5 threshold over
  loaded layer names and their stems.
  \item An \emph{orphan-noun detector} that scans the query for
  noun-like tokens not present in the loaded layer set, filtering GIS
  verbs (buffer, intersect, clip, \ldots) via a stop-word list.
\end{enumerate}
The plugin \emph{never} fetches data from OpenStreetMap. The server
does, but only when the referenced feature name is recognisable
(e.g. \texttt{hospitals}, \texttt{schools}, \texttt{roads}) and the
layer is genuinely absent.

\section{Example Workflows}
\subsection{Proximity + Overlay}
\begin{lstlisting}
User: "Find residential zones within 500m of hospitals"

intent    = PROXIMITY
entities  = {primary: hospitals, secondary: residential}
params    = {distance_value: 500, unit: meters}

plan      = [ buffer(hospitals, 500),
              intersect(step_1_output, residential) ]
\end{lstlisting}

\subsection{Raster NDVI}
\begin{lstlisting}
User: "Calculate NDVI from Sentinel2_red and Sentinel2_NIR"

plan = [ ndvi(RED=Sentinel2_red, NIR=Sentinel2_NIR) ]
\end{lstlisting}

\subsection{Clustering + Hull}
\begin{lstlisting}
User: "Cluster Assets with DBSCAN, then compute the convex hull of each cluster"

plan = [ dbscan(Assets, eps=500),
         convex_hull(step_1_output, group_by=cluster_id) ]
\end{lstlisting}

\section{Public API}
\begin{lstlisting}[language=Python]
from geosi_engine import GeoSI
from geosi_engine.models import Layer

engine = GeoSI()
engine.state.add_layer(Layer(
    name="Schools", layer_type="vector",
    geometry_type="Point", feature_count=1024,
    crs="EPSG:4326", filepath="/data/Schools.shp"
))

result = engine.query("Buffer Schools by 500m")
print(result["success"])       # bool
print(result["answer"])        # user-facing message
print(result["plan"])          # list of ToolStep
print(result["reasoning"])     # LLM or rule-based explanation
\end{lstlisting}

Lower-level access is also available:
\begin{lstlisting}[language=Python]
engine.registry.list_tools(category="vector")
engine.registry.execute_tool("buffer", INPUT="Schools", DISTANCE=500)
\end{lstlisting}

\section{Evaluation}
We exercised GeoSI against the prompt cookbook
(\texttt{prompts.md}) covering inspection, proximity, overlay,
geometry, joins, transformation, sampling, cartography, validation,
statistics, terrain, raster algebra, zonal/focal, resampling,
interpolation, classification, temporal, network, AI/ML, and compound
workflows. Representative outcomes:
\begin{itemize}[leftmargin=*]
  \item Single-step portable queries (e.g. \emph{count features in
  Schools}) run without QGIS, end-to-end.
  \item QGIS-backed queries produce correct \texttt{ExecutionPlan}
  structures; dispatch to QGIS Processing succeeds when the plugin is
  loaded.
  \item Layer-not-found queries consistently produce a useful error
  message with fuzzy suggestions when a near match exists.
  \item All LLM providers exhibit sub-second failover when unavailable.
\end{itemize}

\section{Related Work}
GeoSI builds on several lines of work: natural-language interfaces to
databases (NL2SQL), tool-augmented language models (ReAct, ToolLLM),
and GIS automation frameworks (QGIS Processing, GRASS, WhiteboxTools).
Its novelty lies in the combination: a conversational surface,
strictly grounded in the user's loaded layers, with deterministic
fallback and an explainable plan trace, purpose-built for a
professional GIS environment.

\section{Future Work}
\begin{itemize}[leftmargin=*]
  \item \textbf{v2.5} --- GeoPandas backend for the REST server,
  enabling cloud deployments without QGIS; batch mode.
  \item \textbf{v3.0} --- Real-time satellite stream monitoring, 3D
  viewshed, immersive (XR) visualisation.
  \item \textbf{v4.0} --- Predictive spatial modelling and autonomous
  decision support loops for disaster response and infrastructure
  planning.
\end{itemize}

\section{Conclusion}
GeoSI demonstrates that a small, disciplined architecture can
deliver a conversational GIS that is both grounded and capable. By
keeping the engine framework-free, preferring local LLMs, refusing to
hallucinate data, and always providing a readable failure path, the
system is genuinely deployable today. It is, we believe, the
beginning of geospatial superintelligence --- a domain-native,
explainable, and locally-operable intelligence layer for the GIS
profession.

\vspace{1em}
\noindent\textbf{Availability.} Source code, documentation, and
example datasets are included in the GeoSI repository.

\section*{Acknowledgements}
The author thanks Digital University Kerala (DUK) for institutional
support and the open-source QGIS, Ollama, and scientific Python
communities whose work made GeoSI possible.

\begin{thebibliography}{9}
\bibitem{qgis} QGIS Development Team. \emph{QGIS Geographic Information
System}. Open Source Geospatial Foundation Project, 2024.
\bibitem{ollama} Ollama Project. \emph{Local large language model
runtime}. \url{https://ollama.com}, 2024.
\bibitem{react} Yao, S. et al. \emph{ReAct: Synergizing Reasoning and
Acting in Language Models}. ICLR 2023.
\bibitem{osm} OpenStreetMap contributors. \emph{OpenStreetMap}.
\url{https://www.openstreetmap.org}, 2024.
\end{thebibliography}

\end{document}
